\section{Quadrics from eigenvalue signatures}

Consider
\begin{equation}
 q(\state)=\state\trans\matr A\state+2\vect b\trans\state+c=0,
 \qquad \matr A=\matr A\trans.\label{eq:general-quadric}
\end{equation}
Why introduce eigenvectors?  The original coordinates can be coupled, so
changing one current may change several terms at once.  We ask whether there
is a direction $\vect v$ for which the quadratic map changes only magnitude,
not direction:
\begin{equation}
 \matr A\vect v=\lambda\vect v.
\end{equation}
In a basis made of such directions, every coordinate contributes
independently.  The signs of the corresponding $\lambda$ values then say
whether the surface curves upward, downward, or not at all in each direction.
Accordingly, an orthogonal eigenbasis transforms the quadratic part to
$\lambda_1z_1^2+\lambda_2z_2^2+\lambda_3z_3^2$.  If $\matr A$ is nonsingular,
translation by $-\matr A^{-1}\vect b$ removes all linear terms.  The signs of
the eigenvalues and the translated level $k$ then decide the real surface:
\begin{itemize}
 \item three equal nonzero signs: ellipsoid for the compatible sign of $k$,
       a point for $k=0$, and empty otherwise;
 \item signature $(+,+,-)$: one-sheet hyperboloid for $k>0$, two-sheet
       hyperboloid for $k<0$, and a cone for $k=0$; reversing all signs reverses
       the roles of $k>0$ and $k<0$;
 \item a zero eigenvalue with no linear component in that null direction:
       a cylinder over the corresponding 2D conic;
 \item a zero eigenvalue with a nonzero linear component in that direction:
       an elliptic or hyperbolic paraboloid, according to whether the remaining
       two eigenvalues have equal or opposite signs.
\end{itemize}
Further zero eigenvalues or $k=0$ produce planes, lines, cones, or empty
degeneracies.  This classification follows from principal coordinates; names
are consequences, not assumptions.

The voltage limit is
\begin{equation}
 \norm{\Ae\state}^2\leq\Umax^2
 \quad\Longleftrightarrow\quad
 \state\trans\Qu\state\leq\Umax^2,
 \qquad \Qu=\Ae\trans\Ae.\label{eq:eesm-voltage-quadratic}
\end{equation}
Explicitly,
\begin{equation}
\Qu=\begin{bmatrix}
\Rs^2+\omegae^2\Ld^2&\Rs\omegae(\Ld-\Lq)&\omegae^2\Ld\Lexc\\
\Rs\omegae(\Ld-\Lq)&\Rs^2+\omegae^2\Lq^2&\Rs\omegae\Lexc\\
\omegae^2\Ld\Lexc&\Rs\omegae\Lexc&\omegae^2\Lexc^2
\end{bmatrix}.\label{eq:Qu-explicit}
\end{equation}
It is positive semidefinite because
$\state\trans\Qu\state=\norm{\Ae\state}^2\geq0$, but it cannot be positive
definite: $\rank\Qu=\rank\Ae\leq2<3$.  Under ordinary nondegenerate conditions
$\rank\Ae=2$, hence its eigenvalue signature is $(+,+,0)$.  Since there is no
linear term along the null direction and $U_{\max}^2>0$, the boundary is an
elliptic cylinder.

The determinant makes the same collapse visible as a volume statement.
$\det\Qu=0$ means that the quadratic map cannot scale a three-dimensional
volume to another nonzero volume; one dimension has been flattened.  Here the
reason is structural: two voltage outputs cannot distinguish all directions
in a three-dimensional current input.  The zero determinant, zero eigenvalue,
nontrivial null space, and infinite cylinder axis are therefore not separate
accidents.

\paragraph{Rank-two thought experiment.}
Before returning to machine parameters, consider
\begin{equation}
 \matr A_0=\begin{bmatrix}1&0&1\\0&1&0\end{bmatrix}.
\end{equation}
By inspection, $[-1,0,1]\trans$ is invisible to the map.  The constraint
$\norm{\matr A_0\vect x}^2\leq1$ reads
$(x_1+x_3)^2+x_2^2\leq1$: a unit disk in the two visible combinations and no
restriction along the invisible combination.  This is already an oblique
cylinder.  The EESM matrix changes the directions and scales, not the rank
logic.

\subsection{Structural and parameter-dependent facts}

\textbf{Structural:} because $\Ae\in\R^{2\times3}$,
$\rank(\Ae\trans\Ae)\leq2$.  No positive values of
$R_s,\omega_e,L_d,L_q,L_e$ can make the voltage-only Gram matrix positive
definite in three dimensions.

\textbf{Parameter dependent:} those quantities do determine whether the
ordinary rank is two, the orientation of the null line, and the two finite
semi-axis lengths.  They deform and tilt the cylinder, but they do not close
its axis.

\paragraph{What if there were three independent voltage outputs?}
A full-rank map from $\R^3$ to $\R^3$ could produce a positive-definite Gram
matrix and hence an ellipsoid.  It is the output dimension and rank, not the
word excitation, that forces the cylinder here.
